\scanpage{21}
\mspage{1}
\margin{Nov. Déc. 196\deleted{\unsure{7}}\inserted{6}}
\uline{Cohomologie de De Rham des var. alg.}

\margin{But : faire sentir un complexe de P$^{\mathrm{bs}}$ \unsure{qui existent} autour de la coh. de De Rham des V.A., et proposer un point de vue qui \deleted{\unsure{donne}} me semble permettre de les grouper et de les étudier. D'où un gros programme \unsure{travail proposé}.}
\figure{\begin{tikzpicture}[x=1cm,y=1cm]
\draw (0,0)--(5.1,0)--(5.1,-.40)--(7.8,-.40);
\end{tikzpicture}}

§ 1. \uline{Yoga Cohomologie de De Rham}

\unsure{b)|1)} \uline{Variétés différentiables}. Énoncé du th. de De Rham.
Son explication en termes de théorie des faisceaux :
\deleted{\unsure{i.e.}} $\underline\Omega_X^\bullet$ est une \uline{résolution} de $\mathbb C_X$ (lemme de Poincaré), et les $\underline\Omega_X^i$ sont fins, donc à cohomologie nulle en dim $>0$ (N.B. $X$ est supposé séparé $\ldots$).

2) \uline{Variétés analytiques complexes}. Se \unsure{recopie} tel quel pour les Steins : lemme de Poincaré \inserted{(local)} valable, mais il faut $H^j(X,\underline\Omega_X^i)=0$ pour $j>0$. En général, il faut prendre l'hypercohomologie, on trouve
\diagram{\begin{tikzcd}[ampersand replacement=\&]
H^*(X,\mathbb C) \arrow[d,equal] \&
E_1^{pq}=H^q(X,\underline\Omega_X^p) \arrow[l,Rightarrow] \\
\mathbb H^*(X,\underline\Omega_X^\bullet) \&
\end{tikzcd}}
\margin{*)}

3) \uline{Variétés algébriques} (non singulières). Travaillons avec la top de Zariski. Considérons
\[
\mathbb H^*(X,\underline\Omega_X^\bullet)
\qquad(\text{hypercohomologie})
\]
est-il isomorphe à $H^*(X^{\mathrm{an}},\mathbb C)$ ? On a en tous cas
\diagram{\begin{tikzcd}[ampersand replacement=\&]
\mathbb H^*(X,\underline\Omega_X^\bullet) \arrow[r] \&
\mathbb H^*(X^{\mathrm{an}},\underline\Omega_{X^{\mathrm{an}}})
\simeq H^*(X^{\mathrm{an}},\mathbb C)
\end{tikzcd}}
\margin{20}
